% =============================================================================
% Conclusion (unnumbered, but appears in TOC)
% -----------------------------------------------------------------------------
% A thesis conclusion typically does several jobs:
%   1. Recap the thesis's central question and the path taken to address it.
%   2. Walk through each chapter's contribution, briefly, framing how each
%      built on the previous one.
%   3. Synthesise — what does the thesis say cumulatively that no individual
%      chapter quite says on its own?
%   4. State the contributions explicitly (formal, applied, programmatic).
%   5. Acknowledge limitations.
%   6. Outline directions for future work.
%   7. Discuss implications across the relevant fields.
%   8. Close with a personal / reflective paragraph.
%
% Conclusions are typically 3,000–5,000 words. Use first-person voice
% sparingly but where it earns its place (in synthesis claims and in the
% closing reflection).
% =============================================================================
\chapter*{Conclusion}
\addcontentsline{toc}{chapter}{Conclusion}
\markboth{Conclusion}{Conclusion}
\label{ch:conclusion}

\section*{Where the thesis began}

[Recap the central claim of the thesis and the motivations for taking up
the project.]

\section*{The journey through the chapters}

[One paragraph per chapter, focused on each chapter's contribution to the
through-line of the thesis (not just a summary).]

\section*{Synthesis}

[What does the thesis say taken as a whole? This is where you make claims
that emerge only when the chapters are read together.]

\section*{Original contributions}

[Enumerate the contributions, ideally grouped (e.g., formal, applied,
programmatic). Examiners look for this section.]

\section*{Limitations}

[Be honest about what the thesis does not do. A short paragraph or two.]

\section*{Future work}

[Directions the thesis opens up. Often grouped into 2–4 strands.]

\section*{Implications across fields}

[How the work speaks to the various fields it sits at the intersection of.]

\section*{Closing reflection}

[A short personal close. The conclusion is one of the few places in a
thesis where personal voice is welcome.]
